
\begin{table}[!h]
   \centering
   \small
   \begin{threeparttable}[b]
      \caption{\label{tab:encaje_loan_growth} Reserve-requirement shocks and loan growth: credit lines versus other credit}
      \bigskip
      \begin{tabular}{lcccc}
         \toprule
         Shock measure & \multicolumn{2}{c}{Minimum legal rate} & \multicolumn{2}{c}{Marginal rate} \\ 
         \cmidrule(lr){2-3} \cmidrule(lr){4-5}
         Credit type                       & Credit lines          & Other credit          & Credit lines          & Other credit \\   
                                           & (1)                   & (2)                   & (3)                   & (4)\\  
         \midrule 
         Bank shock                        & -1.055                & 0.082                 & -0.120                & 0.048\\   
                                           & (0.770)               & (0.138)               & (0.113)               & (0.056)\\   
         Banks (= bank clusters)           & 8                     & 11                    & 8                     & 11\\  
         Firms                             & 130                   & 394                   & 95                    & 312\\  
         Firm-by-date cells                & 4,191                 & 17,760                & 1,693                 & 7,187\\  
         Wild cluster bootstrap $p$-value  & 0.560                 & 0.404                 & 0.251                 & 0.565\\  
         Mean of loan growth               & -0.015                & -0.019                & -0.009                & -0.017\\  
          \\
         Observations                      & 9,230                 & 52,029                & 3,686                 & 20,921\\  
         R$^2$                             & 0.482                 & 0.352                 & 0.495                 & 0.368\\  
         Within R$^2$                      & $4.56\times 10^{-5}$  & $2.23\times 10^{-6}$  & $4.16\times 10^{-5}$  & $5.23\times 10^{-5}$\\   
          \\
         Firm $\times$ date fixed effects  & $\checkmark$          & $\checkmark$          & $\checkmark$          & $\checkmark$\\   
         \bottomrule
      \end{tabular}
      
      \begin{tablenotes}\item \textit{Notes:} The dependent variable is the log change in a firm's credit balance (\textit{cr\'edito vigente}) at a bank between consecutive month-ends, $\log B_{bit}-\log B_{bi,t-1}$, so it measures growth among continuing relationships with strictly positive balances at both dates; loan entry and exit are not captured. Balances are summed across products and currencies within a bank--firm--credit-line status (foreign-currency balances are reported in soles at the month-end exchange rate, so growth includes exchange-rate revaluation). Columns (1) and (3) use credit lines (credit cards, current accounts, letters of credit, guarantees, multi-lines, other lines); columns (2) and (4) use all other credit. Because the balance is performing drawn credit, letters of credit, guarantees, multi-lines and other lines have a zero balance in every record, so the credit-line samples consist of credit cards and current accounts. The bank shock is $s^{mn}_{b,0}\Delta r^{mn}_{t}+s^{me}_{b,0}\Delta r^{me}_{t}$, where $\Delta r$ is the month-end change, in percentage points, of the statutory minimum legal reserve rate (columns 1--2) or marginal reserve rate (columns 3--4) on domestic (mn) and foreign (me) currency obligations (positive = tighter), and $s_{b,0}$ are the bank's fixed December 2009 deposit shares in each currency. The marginal-rate shock is undefined when the domestic regime is a single rate, which shortens columns (3)--(4). Each regression keeps firm-by-date cells with at least two banks and absorbs firm-by-date fixed effects. Standard errors, in parentheses, are clustered by bank; with few clusters the $p$-value from a wild cluster bootstrap (null imposed, 9,999 draws) is reported. * $p<0.10$, ** $p<0.05$, *** $p<0.01$ from the conventional clustered $t(G-1)$ test.
      \end{tablenotes}
   \end{threeparttable}
\end{table}


